\begin{proposition}[中央化子中的二次范数与特征 \texorpdfstring{$5$}{5} 的平方像现象]\label{prop:pom-centralizer-det-norm-mod5}
令
\[
M=\begin{pmatrix}1&1\\[2pt]1&0\end{pmatrix},
\qquad
C(M):=\{A\in M_2(\ZZ):\ AM=MA\}.
\]
则 \(C(M)=\ZZ[M]\)，且对任意 \(A\in \ZZ[M]\) 唯一可写为 \(A=aI+bM\)（\(a,b\in\ZZ\)），并满足
\begin{equation}\label{eq:pom-centralizer-det-norm-form}
\det(A)=a^2+ab-b^2.
\end{equation}
此外，在模 \(5\) 意义下恒有
\begin{equation}\label{eq:pom-centralizer-det-square-mod5}
\det(A)\equiv (a+3b)^2\pmod 5,
\end{equation}
从而 \(\det(\ZZ[M]\otimes \FF_5)\) 的像包含于 \(\FF_5\) 的平方元集合。
反之，对任意素数 \(p\neq 5\)，\(\det(\ZZ[M]\otimes \FF_p)\) 的像为整个 \(\FF_p\)。
\end{proposition}
\begin{proof}
由定理 \ref{thm:pom-centralizer-semiring-fibonacci-fusion} 的整数版本，\(C(M)=\ZZ[M]\) 且任意 \(A\in\ZZ[M]\) 唯一写作 \(A=aI+bM\)。
直接计算
\[
A=aI+bM
=\begin{pmatrix}a+b&b\\[2pt]b&a\end{pmatrix},
\qquad
\det(A)=(a+b)a-b^2=a^2+ab-b^2,
\]
得到 \eqref{eq:pom-centralizer-det-norm-form}。
\par
在 \(\FF_5\) 中有 \(-1\equiv 4\)，故
\(
a^2+ab-b^2\equiv a^2+ab+4b^2
\)。
取 \(k\equiv 3\pmod 5\)，则 \(2k\equiv 1\pmod 5\) 且 \(k^2\equiv 9\equiv 4\pmod 5\)，于是
\[
(a+kb)^2\equiv a^2+2kab+k^2b^2\equiv a^2+ab+4b^2\pmod 5,
\]
即得 \eqref{eq:pom-centralizer-det-square-mod5}。
\par
对 \(p\neq 5\)，二次多项式 \(\tau^2-\tau-1\) 的判别式为 \(5\)，在 \(\FF_p\) 上不为零，故
\(
\FF_p[\tau]/(\tau^2-\tau-1)
\)
为二次 \'{e}tale \(\FF_p\)-代数（分裂为 \(\FF_p\times\FF_p\) 或为二次扩张 \(\FF_{p^2}\)）。
在同构 \(\ZZ[M]\otimes\FF_p\cong \FF_p[\tau]/(\tau^2-\tau-1)\) 下，\(\det(aI+bM)\) 恰为该二次 \'{e}tale 代数到 \(\FF_p\) 的范数。
而有限域扩张的范数映射在乘法群上满射，分裂情形亦为两分量乘积之满射；故 \(\det\) 的像为 \(\FF_p\) 的全体元素。
\leanverified{centralizer\_det\_formula}
\leanverified{det\_mod5\_is\_square}
\leanverified{golden\_ratio\_discriminant}
\leanverified{golden\_poly\_mod5\_double\_root}
\leanverified{paper\_pom\_centralizer\_det\_mod5}
\leanverified{paper\_pom\_centralizer\_det\_norm\_mod5}
\end{proof}

\begin{corollary}[判别式锁定的唯一坏素数：\texorpdfstring{$p=5$}{p=5} 的骨架退化]\label{cor:pom-centralizer-unique-bad-prime-5}
在中央化子骨架 \(\ZZ[M]\cong \ZZ[\tau]/(\tau^2-\tau-1)\) 上，模素约化 \(\ZZ[M]\otimes\FF_p\) 为非 \'{e}tale（非约化）当且仅当 \(p=5\)；
等价地，特征 \(5\) 是该二次骨架在纯判别式层面发生退化的唯一算术奇点。
\end{corollary}
\begin{proof}
\(\tau^2-\tau-1\) 的判别式为 \(5\)。
对 \(p\neq 5\)，\(\FF_p[\tau]/(\tau^2-\tau-1)\) 为 \'{e}tale 代数；
对 \(p=5\)，有 \(\tau^2-\tau-1\equiv(\tau-3)^2\) 于 \(\FF_5\)，从而约化环非约化并带幂零元。
由 \(\ZZ[M]\cong \ZZ[\tau]/(\tau^2-\tau-1)\) 得证。
\end{proof}
\leanverified{paper\_pom\_centralizer\_unique\_bad\_prime\_5}
